\DTMsavedate{mydate}{2022-06-30}
\maketitle{Network Theory}{Electronic Engineering}{\DTMusedate{mydate}}

% ----------------------------------------------------- %
% COURSE INFO
\subsection*{Course Information}
\vspace{0ex}%
\noindent\parbox{0.5\textwidth}{%
    \noindent\begin{tabular}{@{}ll}
                 \textsf{Course:}          & 	Network Theory \\
                 \textsf{Course Code:}         & 	TEE 2101    \\
                 \textsf{Credit Hours:}    & 	10
    \end{tabular}}
\vspace{2ex}\hrule\vspace{2ex}

% ----------------------------------------------------- %
% INSTRUCTOR INFO
\subsection*{Lecturer Information}
\noindent\parbox{0.5\textwidth}{%
    \noindent\begin{tabular}{@{}ll}
                 \textsf{Name:}         &     Magripa Nleya         \\
%\textsf{Phone:}        &    \phone          \\
\textsf{Email:}        &    magripa.nleya@nust.ac.zw \\
\textsf{Qualifications:}        &  MSc in Telecommunication Engineering
    \end{tabular}}
\vspace{2ex}\hrule\vspace{2ex}


% ----------------------------------------------------- %
% COURSE DESCRIPTION
\subsection*{Learning Aim}
To instill concepts of Network Theory into students.
\vspace{2ex} \hrule\vspace{2ex}

% ----------------------------------------------------- %
% LEARNING OBJECTIVES
\subsection*{Learning Objectives}
By the end of the course, candidates should be able to:
\begin{outline}
    \1 analyse electrical circuits.

    \1 identify the type of the circuit and solve problems using appropriate method for the given circuit.
\end{outline}
\vspace{2ex}\hrule\vspace{2ex}

% ----------------------------------------------------- %
% COURSE TOPICS
\subsection*{Topics}
\begin{outline}
    \1 DC circuits analysis
    \begin{enumerate}
        \item First order circuits
        \item Second order circuits
    \end{enumerate}
    \1 Analysing AC circuits analysis
    \begin{enumerate}
        \item Sinusoids and phasors
        \item Sinusoidal steady-state analysis
        \item Frequency response
    \end{enumerate}
    \1 Advanced circuit analysis
    \begin{enumerate}
        \item The Laplace Transform
        \item The Fourier series
        \item Fourier Transform
        \item Two-port networks
    \end{enumerate}
\end{outline}
\vspace{2ex} \hrule\vspace{2ex}

% ----------------------------------------------------- %
% Students Assenssment
\subsection*{Assessment of Students}
Tests, Laboratory, Assignments and Final Examination at the end of the semester.
\vspace{2ex}\hrule\vspace{2ex}

% ----------------------------------------------------- %
% Grade Evaluation
\subsection*{Course Evaluation and Grading Policies}
Grades are based on:
Continuous assessment (\qty{25}{\percent}),
Final Exam (\qty{75}{\percent})
\vspace{2ex}\hrule\vspace{2ex}
% ----------------------------------------------------- %
% EQUIPMENT
% https://sites.udel.edu/computing-purchases/personal-specs/
% \subsection*{Essential Equipment}

% \vspace{2ex}\hrule\vspace{2ex}

% Policies and Other information
\subsection*{Policies and Other Information}
% Key Text or some Books
\subsection*{Key Text}

\begin{itemize}
    \item Introductory circuit analysis by Boylestard
    \item Fundamentals of electric circuits by C.K. Alexander, M.N.O. Sadiku.
    \item Lecture notes.
\end{itemize}
\vspace{2ex}\hrule
\vspace{2ex}\hrule\vspace{2ex}